\chapter*{Abstract}
\addcontentsline{toc}{chapter}{Abstract}

The digital transformation of the labor market has led to an abundance of online internship listings, yet this "data deluge" presents significant challenges due to information being unstructured, inconsistent, and highly fragmented across diverse platforms. Traditional recruitment systems often rely on rigid lexical matching, which fails to bridge the semantic gap between varying job titles and skill descriptions, while modern Large Language Model (LLM) approaches remain computationally expensive and latent for large-scale real-time ingestion. This thesis proposes a lightweight, tiered semantic enrichment pipeline designed to transform chaotic internship data into a structured, searchable knowledge base without the overhead of heavy computational infrastructure.

Our approach introduces a modular architecture that combines deterministic rule-based extraction, lightweight neural mapping using transformer models (such as DistilBERT) for standardization against the ESCO professional taxonomy, and a resource-budgeted LLM fallback for complex cases. Furthermore, we develop a flexible ingestion system capable of automatically normalizing multi-format JSON data through fuzzy-key matching and content-based inference. The processed data is stored in a hybrid retrieval system utilizing PostgreSQL and the \texttt{pgvector} extension, enabling sub-50ms response times by blending lexical (BM25) and semantic (HNSW) search.

The system was evaluated using a real-world dataset of 17,500 listings. Results demonstrate that the hybrid search strategy significantly improves relevance and recall compared to traditional keyword-based methods. Moreover, the tiered pipeline achieves high semantic accuracy while maintaining production-ready efficiency on commodity hardware. This research provides a scalable and cost-effective framework for universities and recruitment platforms to manage and retrieve internship opportunities with high precision.
